\chapter*{Lampiran}
\addcontentsline{toc}{chapter}{Lampiran}

% ============================================================
% Lampiran A: Rubrik Penilaian
% ============================================================
\section*{Lampiran A: Rubrik Penilaian Tugas Praktik Web}

Rubrik berikut terhubung dengan CPMK dan Sub-CPMK RPS Pemrograman Web. Setiap kriteria memetakan ke capaian pembelajaran yang dapat diukur.

\begin{table}[htbp]
\centering
\small
\begin{tabular}{|>{\raggedright\arraybackslash}p{2.6cm}|>{\raggedright\arraybackslash}p{2.2cm}|>{\raggedright\arraybackslash}p{3cm}|>{\raggedright\arraybackslash}p{3cm}|>{\raggedright\arraybackslash}p{3cm}|}
\hline
\textbf{Kriteria} & \textbf{Sub-CPMK} & \textbf{Sangat Baik} & \textbf{Baik} & \textbf{Perlu Perbaikan} \\
\hline
Struktur HTML & 2.1 & Semantik tepat, aksesibel & Semantik cukup baik & Banyak elemen tidak tepat \\
\hline
CSS dan Layout & 2.2 & Responsif, rapi, konsisten & Cukup responsif & Layout berantakan \\
\hline
Interaktivitas JS & 2.3 & Validasi dan event tepat & Ada fitur utama & Banyak error/interaksi minim \\
\hline
Back-end & 3.1, 3.2 & CRUD lengkap, aman & CRUD sebagian & Banyak error/fungsi hilang \\
\hline
Keamanan Dasar & 4.2 & Validasi input dan sanitasi & Validasi sebagian & Tidak ada validasi \\
\hline
\end{tabular}
\caption{Rubrik Penilaian Tugas Praktik Web (terhubung CPMK/Sub-CPMK)}
\end{table}

\textbf{Pemetaan ke CPMK:} Struktur HTML $\to$ Sub-CPMK 2.1 (HTML5); CSS dan Layout $\to$ Sub-CPMK 2.2 (CSS3, responsif); Interaktivitas JS $\to$ Sub-CPMK 2.3 (DOM, event, validasi form); Back-end $\to$ Sub-CPMK 3.1 (PHP), 3.2 (MySQL, CRUD); Keamanan Dasar $\to$ Sub-CPMK 4.2 (validasi, sanitasi, XSS/CSRF). Asesmen UTS mengacu CPMK-1 dan CPMK-2; UAS dan proyek mengacu CPMK-1 s.d. CPMK-4 sesuai bobot RPS.

% ============================================================
% Lampiran B: Contoh Template Laporan
% ============================================================
\section*{Lampiran B: Contoh Template Laporan Tugas}
\begin{enumerate}
  \item Judul dan Identitas Mahasiswa
  \item Deskripsi Masalah dan Kebutuhan Pengguna
  \item Desain UI/UX (wireframe dan mockup singkat)
  \item Implementasi (cuplikan kode penting)
  \item Hasil Pengujian dan Evaluasi
  \item Refleksi dan Kesimpulan
\end{enumerate}

% ============================================================
% Lampiran C: Glosarium
% ============================================================
\section*{Lampiran C: Glosarium Istilah Web}
\begin{itemize}
  \item \textbf{HTML}: Bahasa markup untuk struktur halaman web.
  \item \textbf{CSS}: Bahasa untuk styling dan layout halaman web.
  \item \textbf{DOM}: Representasi struktur dokumen yang dapat dimanipulasi.
  \item \textbf{API}: Antarmuka yang memungkinkan pertukaran data antar sistem.
  \item \textbf{CRUD}: Operasi Create, Read, Update, Delete pada data.
  \item \textbf{Responsive}: Desain yang menyesuaikan ukuran layar.
  \item \textbf{Session}: Mekanisme menyimpan informasi pengguna selama interaksi.
  \item \textbf{XSS/CSRF}: Jenis serangan keamanan pada aplikasi web.
\end{itemize}

% ============================================================
% Lampiran D: Referensi Tambahan
% ============================================================
\section*{Lampiran D: Referensi Tambahan}
\begin{itemize}
  \item MDN Web Docs: \url{https://developer.mozilla.org}
  \item W3Schools: \url{https://www.w3schools.com}
  \item OWASP Cheat Sheets: \url{https://cheatsheetseries.owasp.org}
\end{itemize}

